\twocolumn[
\maketitle

\selectlanguage{english}
\begin{abstract}
Municipal governments frequently need to draft or adapt Bills under severe time and evidence constraints. This paper proposes and evaluates a reproducible pipeline to discover Brazilian SAPL instances, extract municipal Bills, translate legislative summaries into operational actions, retrieve semantically related actions, and associate recurrent proposals with official municipal indicators. In a frozen execution dated November 12, 2025, the pipeline discovered 1,259 SAPL instances, achieved automated extraction in 322 of them, and collected 220,065 Bills. The text-to-action translator was trained on 1,000 synthetic pairs derived from real summaries and obtained a BERTScore of 84\% on the validation split. The resulting system, Sonar Municipal, operationalizes the pipeline in a web platform for non-technical users. We report the current methodological scope, preliminary results, reproducibility artifacts, and the remaining threats to validity, especially the absence of a formal retrieval benchmark and the non-causal nature of the indicator analysis.
\keywords{municipal legislation; SAPL; semantic search; policy recommendation; open data; reproducibility}
\end{abstract}

\selectlanguage{brazil}
\begin{resumo}
Gestores municipais frequentemente precisam redigir, adaptar ou comparar Projetos de Lei (PLs) com forte restrição de tempo e de evidências. Este artigo propõe e avalia um pipeline reprodutível para descobrir instâncias brasileiras do SAPL, extrair PLs municipais, traduzir ementas em ações operacionais, recuperar ações semanticamente relacionadas e associar propostas recorrentes a indicadores oficiais municipais. Em uma execução congelada em 12 de novembro de 2025, o pipeline encontrou 1.259 instâncias do SAPL, obteve extração automatizada em 322 delas e coletou 220.065 PLs. O tradutor \textit{ementa} $\rightarrow$ \textit{ação} foi treinado com 1.000 pares sintéticos derivados de ementas reais e atingiu BERTScore de 84\% no conjunto de validação. O sistema resultante, Sonar Municipal, operacionaliza o pipeline em uma plataforma web para usuários não técnicos. Apresentamos o escopo metodológico atual, resultados preliminares, artefatos de reprodutibilidade e as ameaças à validade ainda em aberto, em especial a ausência de um benchmark formal de recuperação e o caráter não causal da análise com indicadores.
\palavrasChave{legislação municipal; SAPL; busca semântica; recomendação de políticas; dados abertos; reprodutibilidade}
\end{resumo}
\vspace{12pt}
]
